\section{Introduction}

Outcome-based reinforcement learning has become a standard post-training tool for improving the reasoning ability of language models. In mathematical reasoning, a verifier can often judge whether a completed response is correct, and policy optimization can then increase the probability of successful responses while decreasing the probability of failed ones \cite{schulman2017proximal,ouyang2022training,shao2024deepseekmath,guo2025deepseekr1}. Group-relative methods such as GRPO make this setup especially simple: several responses are sampled for the same prompt, a scalar outcome reward is computed for each response, and the resulting relative advantage is broadcast to all response tokens. This design uses reliable final feedback, but it provides little information about which intermediate tokens should be reinforced.

On-policy distillation (OPD) offers a complementary source of supervision. Instead of training only from a scalar outcome, the student first samples its own response and a stronger teacher is then evaluated on the prefixes actually visited by the student. The teacher provides dense token-level guidance, often through the log-probability gap between the teacher and the student. This keeps training on-policy while exposing the student to fine-grained teacher preferences. For long reasoning traces, however, dense teacher supervision is not uniformly reliable. The teacher may locally prefer a token that belongs to an ultimately incorrect solution, or locally disfavor a token that appears in a correct trajectory. This mismatch creates a credit-assignment problem: the local teacher signal and the final outcome do not always agree.

A second issue arises from the prefixes on which the teacher is queried. OPD evaluates the teacher on student-generated prefixes rather than on teacher-generated ones. Once a student trajectory drifts into a low-support region under the teacher, later teacher logit differences still exist, but they may be less informative about how to reach a correct solution. Treating these later dense rewards as uniformly trustworthy can amplify supervision from invalid or off-path prefixes. Thus, OPD needs not only token-outcome consistency, but also a way to judge whether the current prefix remains a valid context for teacher supervision.

We propose \emph{Prefix-Outcome Validation} (POV), a simple validation mechanism for on-policy distillation. POV uses final correctness to validate the direction of dense teacher preferences, rather than replacing them with a sparse reward. In a correct trajectory, a candidate token is outcome-aligned when the teacher assigns it higher log-probability than the student. In an incorrect trajectory, the alignment direction is reversed: a candidate is outcome-aligned when the teacher assigns it lower log-probability than the student, thereby discouraging behavior observed in a failed response. This rule converts the final verifier result into a token-level reliability test for OPD supervision.

POV also validates prefixes. Since the teacher is queried on student-generated prefixes, the later parts of a response may drift into regions where the teacher's dense preferences are less trustworthy. POV monitors teacher excess surprisal, defined as the negative teacher log-probability of the sampled token minus the teacher entropy at the same prefix. A sustained increase in excess surprisal indicates that the sampled prefix has become unexpectedly unlikely under the teacher. POV detects this change with a windowed CUSUM statistic and monotonically decays subsequent dense supervision after the detected horizon.

The resulting method, POV-OPD, keeps the useful part of dense OPD while reducing the influence of locally unreliable signals. Outcome validation assigns full weight to outcome-aligned token candidates and uses a beta-capped rank weight for anti-aligned candidates. Prefix validation then attenuates supervision after sustained teacher-surprisal drift. Thus, POV-OPD preserves token-level teacher guidance when it is consistent with both the final outcome and the teacher's prefix-level support.

Our contributions are:
\begin{itemize}
    \item We identify two reliability issues in on-policy distillation for reasoning: outcome-token mismatch and prefix drift under student-generated trajectories.
    \item We introduce outcome validation, which uses final correctness to partition dense teacher preferences into outcome-aligned and anti-aligned token candidates.
    \item We introduce prefix validation, which detects sustained excess surprisal under the teacher and monotonically down-weights later dense supervision.
    \item We combine both mechanisms into POV-OPD, a conservative modification of dense OPD that preserves aligned supervision and assigns only beta-capped rank weight to anti-aligned candidates.
\end{itemize}
